\chapter{Całka Itô: losowe funkcje podcałkowe}
\label{ch:ito-construction}
\section{Dlaczego funkcja podcałkowa musi znać tylko przeszłość?}
W rozdziale pierwszym współczynnik $a_k$ w sumie
$\sum a_k\Delta_kW$ był ustaloną liczbą. Jeżeli chcemy
modelować strategię inwestycyjną, siłę losowego zakłócenia
albo przyrost procesu zależny od aktualnego położenia,
to naturalny współczynnik ma postać $H_{t_k}(\omega)$.
Aby rachunek nie wykorzystywał przyszłości, musimy żądać,
by $H_{t_k}$ był mierzalny względem $\calF_{t_k}$.

W przeciwnym razie suma mogłaby znać znak następnego
przyrostu i osiągać dodatnią średnią nawet wtedy, gdy
podstawowy proces Wienera jest martyngałem. Na przykład
$H_{t_k}=\operatorname{sgn}(\Delta_kW)$
uzależnia decyzję od przyrostu, który jeszcze nie nastąpił:
wtedy każdy składnik $H_{t_k}\Delta_kW=|\Delta_kW|$
jest nieujemny. Taka strategia nie jest przewidywalna.

\begin{definicja}[Proces prosty przewidywalny]
Proces $H:[0,T]\times\Omega\to\R$ nazywamy prostym
przewidywalnym, jeżeli istnieje deterministyczny podział
$0=t_0<\cdots<t_n=T$ i zmienne losowe $\xi_k$
mierzalne względem $\calF_{t_k}$, takie że
\[
H_t=\sum_{k=0}^{n-1}\xi_k\ind{(t_k,t_{k+1}]}(t),
\qquad 0\le t\le T.
\tag{2.1}\label{eq:ito-step}
\]
W przestrzeni $\calS_{\mathrm{pred}}^2$ rozpatrujemy
te procesy proste, dla których $\E\xi_k^2<\infty$
na każdym przedziale o dodatniej długości.
\end{definicja}

\begin{uwaga}[Wspólna filtracja ma znaczenie]
Jeżeli do filtracji dołączylibyśmy już w chwili zerowej
wartość $W_T$, to $W$ nie byłby procesem Wienera
\emph{względem powiększonej filtracji}: przyrost
$W_T-W_0$ nie byłby niezależny od $\calF_0$.
Dlatego każda definicja całki stochastycznej zawiera
zarówno proces Wienera, jak i filtrację, względem której
jest on procesem Wienera.
\end{uwaga}

\begin{center}
\begin{tikzpicture}[>=Latex,x=1cm,y=0.9cm]
\draw[->] (0,0)--(9.2,0) node[right] {$t$};
\draw[->] (0,-.35)--(0,2.8) node[above] {$H_t$};
\fill[blue!12] (.2,0) rectangle (2,1.35);
\fill[blue!20] (2,0) rectangle (4.15,2.15);
\fill[blue!12] (4.15,0) rectangle (6.4,.8);
\fill[blue!20] (6.4,0) rectangle (8.45,1.8);
\draw[thick,blue!65!black] (.2,1.35)--(2,1.35) (2,2.15)--(4.15,2.15)
(4.15,.8)--(6.4,.8) (6.4,1.8)--(8.45,1.8);
\foreach \x/\name in {0.2/{t_0},2/{t_1},4.15/{t_2},6.4/{t_3},8.45/{t_4}}{
 \draw[dashed,gray] (\x,0)--(\x,2.5);
 \node[below,scale=.83] at (\x,0) {$\name$};}
\node[scale=.86,blue!55!black] at (1.05,1.68) {$\xi_0$};
\node[scale=.86,blue!55!black] at (3.05,2.48) {$\xi_1$};
\node[scale=.86,blue!55!black] at (5.3,1.13) {$\xi_2$};
\node[scale=.86,blue!55!black] at (7.5,2.12) {$\xi_3$};
\end{tikzpicture}
\par\smallskip\small Przewidywalny proces prosty: wysokość przedziału jest znana przy jego lewym końcu.
\end{center}

\section{Całka prostego procesu i warunkowa izometria}
\begin{definicja}[Całka Itô dla procesów prostych]
Dla $H$ postaci \eqref{eq:ito-step} i $t\in[0,T]$ kładziemy
\[
I_t(H):=\int_0^t H_s\,\dd W_s
:=\sum_{k=0}^{n-1}\xi_k
\big(W_{t\wedge t_{k+1}}-W_{t\wedge t_k}\big).
\tag{2.2}\label{eq:ito-simple}
\]
Zmienna $I_t(H)$ jest $\calF_t$-mierzalna.
\end{definicja}

\begin{lemat}[Poprawność definicji]
Wzór \eqref{eq:ito-simple} nie zależy od wyboru
reprezentacji $H$ jako procesu prostego (z dokładnością
do równości $\dd t\otimes\dd\Prob$-prawie wszędzie).
Dla prostych procesów $H,K$ i $a,b\in\R$ zachodzi
$I_t(aH+bK)=aI_t(H)+bI_t(K)$ niemal na pewno.
\end{lemat}
\begin{proof}
Przy rozdrobnieniu deterministycznego podziału współczynnik
$\xi_k$ na każdym nowym przedziale pozostaje mierzalny
względem $\calF$ jego lewego końca, bo filtracja jest rosnąca.
Przyrost procesu Wienera dzieli się na sumę przyrostów,
więc wartość całki nie zmienia się. Dwie reprezentacje
procesu sprowadzamy do wspólnego podziału. Na każdym
przedziale ich współczynniki są równe niemal na pewno,
ponieważ norma różnicy w przestrzeni produktowej jest zerowa:
$\E\int_0^T|H-K|^2\,\dd t=0$.
Po usunięciu skończenie wielu zbiorów zerowych
współczynniki są równe jednocześnie. Stąd niezależność
od reprezentacji i liniowość.
\end{proof}

\begin{twierdzenie}[Martyngałowość i izometria Itô dla procesów prostych]
\label{th:ito-simple-iso}
Niech $H,K\in\calS_{\mathrm{pred}}^2$. Wtedy
$I(H)=(I_t(H))_{0\le t\le T}$ jest ciągłym
martyngałem z $I_0(H)=0$ i zachodzą
\[
\E I_t(H)=0,\qquad
\E[I_t(H)I_t(K)]
=\E\int_0^tH_sK_s\,\dd s,
\tag{2.3}\label{eq:ito-iso-simple}
\]
\[
\E\big[(I_t(H)-I_s(H))^2\mid\calF_s\big]
=\E\left[\left.\int_s^tH_u^2\,\dd u\,\right|\calF_s\right]
\quad(0\le s\le t\le T).
\tag{2.4}\label{eq:ito-conditional-var}
\]
\end{twierdzenie}
\begin{proof}
Ustalmy wspólny podział dla $H,K$ i początkowo przyjmijmy
$t=T$. Zapiszmy $H=\sum\xi_k\ind{(t_k,t_{k+1}]}$,
$K=\sum\eta_k\ind{(t_k,t_{k+1}]}$ i
$\Delta_kW=W_{t_{k+1}}-W_{t_k}$.
Dla każdej $k$ mamy
$\E[\Delta_kW\mid\calF_{t_k}]=0$ i
$\E[(\Delta_kW)^2\mid\calF_{t_k}]=\Delta t_k$,
ponieważ przyrost jest niezależny od $\calF_{t_k}$.
Stąd
$\E[\xi_k\Delta_kW\mid\calF_{t_k}]=0$ i
$\E[\xi_k\eta_k(\Delta_kW)^2]=\E[\xi_k\eta_k]\Delta t_k$.

Dla $j<k$ zmienna $\xi_j\eta_k\Delta_jW$ jest
$\calF_{t_k}$-mierzalna, zatem
\[
\E[\xi_j\eta_k\Delta_jW\Delta_kW]
=\E\big[\xi_j\eta_k\Delta_jW\,
\E(\Delta_kW\mid\calF_{t_k})\big]=0.
\]
Identycznie znikają wyrazy o $k<j$.
Po wymnożeniu iloczynu dwóch sum pozostają więc
wyłącznie wyrazy diagonalne:
\[
\E[I_T(H)I_T(K)]
=\sum_k\E(\xi_k\eta_k)\Delta t_k
=\E\int_0^TH_sK_s\,\dd s.
\]
Gdy $t<T$, wystarczy uciąć procesy na $(0,t]$,
co dowodzi \eqref{eq:ito-iso-simple}. Analogicznie,
aby pokazać własność martyngałową, dzielimy $I_t(H)$
na $I_s(H)$ i sumę przyrostów z $(s,t]$. Każdy taki
składnik ma warunkową średnią zero względem informacji
w swoim lewym końcu; iterowane warunkowanie daje
$\E[I_t(H)-I_s(H)\mid\calF_s]=0$. Ciągłość trajektorii
wynika ze skończonej sumy ciągłych procesów przyrostowych.

Na koniec zastosujmy identyczne rozwinięcie kwadratu
do całki na przedziale $(s,t]$, ale tym razem z warunkowaniem
względem $\calF_s$. Wyrazy mieszane znikają po iterowanym
warunkowaniu, a składniki diagonalne mają warunkową
wartość oczekiwaną równą
$\E[\xi_k^2\Delta t_k\mid\calF_s]$.
Zsumowanie daje \eqref{eq:ito-conditional-var}.
\end{proof}

\begin{uwaga}[Uzasadnienie całkowalności iloczynów]
Jeżeli $\E\xi_k^2$ i $\E\eta_k^2$ są skończone,
to $\E|\xi_k\eta_k|<\infty$ z nierówności
Cauchy'ego--Schwarza. Same składniki całek mają skończone
drugie momenty dzięki warunkowej wariancji przyszłego
przyrostu. W konsekwencji wszystkie zastosowania
warunkowej wartości oczekiwanej w dowodzie są uprawnione.
\end{uwaga}

\section{Procesy przewidywalne i przestrzeń całkowania}
W praktyce proces $H$ nie musi mieć skończonej liczby zmian.
Na przykład $H_t=W_t$ jest ciągły i zależy od całej
przeszłości. Aby przybliżać takie procesy przez
$\calS_{\mathrm{pred}}^2$, potrzebujemy dokładnie określonej
mierzalności względem czasu i zdarzeń losowych.

\begin{definicja}[Przewidywalne $\sigma$-ciało]
Na przestrzeni $[0,T]\times\Omega$ definiujemy
\emph{przewidywalne $\sigma$-ciało} $\calP_T$ jako
$\sigma$-ciało generowane przez zbiory
$\{0\}\times A$, $A\in\calF_0$, oraz
$(s,t]\times A$, gdzie $0\le s<t\le T$ i $A\in\calF_s$.
Proces $H$ nazywamy \emph{przewidywalnym}, gdy odwzorowanie
$(t,\omega)\mapsto H_t(\omega)$ jest
$\calP_T/\calB(\R)$-mierzalne.
Definiujemy
\[
\calH_T^2:=\left\{H\text{ przewidywalny}:
\E\int_0^T H_t^2\,\dd t<\infty\right\}\big/\sim,
\quad H\sim K\iff H=K\quad
(\dd t\otimes\Prob)\text{-p.w.}
\tag{2.5}
\]
Norma ma postać
$\norm{H}_{\calH_T^2}:=(\E\int_0^TH_t^2\,\dd t)^{1/2}$.
\end{definicja}

\begin{uwaga}[Oznaczenia i mierzalność]
Produktem miar w \eqref{eq:ito-iso-simple} jest
$\dd t\otimes\dd\Prob$, w zwykłym znaczeniu miary
Lebesgue'a na czasie i miary probabilistycznej na $\Omega$.
Nie każdy proces adaptowany jest automatycznie przewidywalny.
Natomiast każdy adaptowany proces o ciągłych albo
lewostronnie ciągłych trajektoriach jest przewidywalny:
można przybliżać go punktowo procesami stałymi na
przedziałach, używając wartości w lewych końcach.
\end{uwaga}

\begin{lemat}[Gęstość prostych procesów przewidywalnych]
\label{lem:pred-dense}
Przestrzeń $\calS_{\mathrm{pred}}^2$ jest gęsta w
$\calH_T^2$.
\end{lemat}
\begin{proof}
Przestrzeń $\calH_T^2$ jest zamkniętą podprzestrzenią
$L^2([0,T]\times\Omega,\calB([0,T])\otimes\calF,
\dd t\otimes\Prob)$ z warunkiem mierzalności względem
$\calP_T$, a więc jest zupełna.
Rozważmy rodzinę prostokątów $(s,t]\times A$ z
$A\in\calF_s$. Ich wskaźniki odpowiadają dokładnie
prostym procesom przewidywalnym
$\ind A\ind{(s,t]}$. Rodzina tych prostokątów, wraz
ze zbiorami z chwili zerowej, jest zamknięta na
skończone przecięcia: jeśli $s\le u$, to
$\big((s,t]\times A\big)\cap\big((u,v]\times B\big)$
jest albo pusta, albo równa
$(u,\min\{t,v\}]\times(A\cap B)$, gdzie
$A\cap B\in\calF_u$. Tak więc jest to $\pi$-układ
generujący $\calP_T$.

Niech $\calD$ będzie klasą zbiorów $C\in\calP_T$,
dla których wskaźnik $\ind C$ należy do domknięcia
w $L^2(\dd t\otimes\Prob)$ przestrzeni liniowej
rozpiętej przez wskaźniki wspomnianych prostokątów.
Ponieważ miara produktowa ma skończoną masę $T$,
$\calD$ jest klasą Dynkina: zawiera całą przestrzeń
(z dokładnością do przekroju $\{0\}\times\Omega$ o
zerowej mierze produktowej), jest zamknięta na różnice
zbiorów zagnieżdżonych oraz na przeliczalne rozłączne
sumy, co wynika ze zbieżności wskaźników w $L^2$
i z domknięcia. Twierdzenie o $\pi$- i $\lambda$-układach
daje $\calD=\calP_T$. Wobec tego wskaźniki dowolnych
zbiorów przewidywalnych dają się aproksymować w $L^2$
przez skończone liniowe kombinacje wskaźników prostokątów.
Na koniec każdą funkcję z $L^2(\calP_T)$ przybliżamy
funkcjami prostymi. Otrzymane aproksymacje są
prostymi procesami przewidywalnymi o ograniczonych
współczynnikach, a zatem należą do
$\calS_{\mathrm{pred}}^2$.
\end{proof}

\begin{twierdzenie}[Konstrukcja całki Itô w $L^2$]\label{th:ito-extension}
Dla każdego $H\in\calH_T^2$ istnieje jednoznaczny
(z dokładnością do nierozróżnialności) ciągły,
adaptowany martyngał $I(H)=(I_t(H))_{0\le t\le T}$
z $I_0(H)=0$, taki że dla każdego ciągu
$H^n\in\calS_{\mathrm{pred}}^2$ z
$\norm{H^n-H}_{\calH_T^2}\to0$ zachodzi
\[
\E\sup_{0\le t\le T}|I_t(H^n)-I_t(H)|^2\longrightarrow0.
\tag{2.6}\label{eq:ito-process-limit}
\]
Odwzorowanie $H\mapsto I_T(H)$ jest liniową izometrią
$\calH_T^2\to L^2(\Omega)$ oraz dla dowolnych
$H,K\in\calH_T^2$ i $t\in[0,T]$ zachodzi
\[
\E\left|\int_0^tH_s\,\dd W_s\right|^2
=\E\int_0^tH_s^2\,\dd s,
\qquad
\E\left[\int_0^tH_s\,\dd W_s\int_0^tK_s\,\dd W_s\right]
=\E\int_0^tH_sK_s\,\dd s.
\tag{2.7}\label{eq:ito-full-iso}
\]
\end{twierdzenie}
\begin{proof}
Gęstość z lematu \ref{lem:pred-dense} zapewnia istnienie
$H^n\to H$ w $\calH_T^2$. Dla procesów prostych
$M^n:=I(H^n)$ nierówność maksymalna Dooba i
\eqref{eq:ito-iso-simple} dają
\[
\E\sup_{t\le T}|M_t^n-M_t^m|^2
\le4\E|M_T^n-M_T^m|^2
=4\norm{H^n-H^m}_{\calH_T^2}^2.
\tag{2.8}\label{eq:ito-doob-cauchy}
\]
W szczególności $(M^n)$ jest ciągiem Cauchy'ego
w normie $\norm{M}_{\mathbb S^2}:=(\E\sup_{t\le T}|M_t|^2)^{1/2}$.
Aby nie opierać się niejawnie na zupełności przestrzeni
procesów, wybierzmy podciąg $(M^{n_j})$, dla którego
$\sum_j\norm{M^{n_{j+1}}-M^{n_j}}_{\mathbb S^2}<\infty$.
Z nierówności Jensena i zbieżności monotonicznej
$\E\sum_j\sup_t|M_t^{n_{j+1}}-M_t^{n_j}|<\infty$.
Zatem suma supremów jest skończona prawie na pewno,
a podciąg zbiega jednostajnie na $[0,T]$ poza jednym
zbiorem zerowym. Jego granica ma ciągłe trajektorie.
Niech $M$ oznacza ten proces, z wartością zero na
wyjątkowym zbiorze (który można wybrać w $\calF_0$
przez zupełność filtracji).

Ponieważ procesy $M^{n_j}$ są adaptowane, ich granice
punktowe są $\calF_t$-mierzalne. Z nierówności
\eqref{eq:ito-doob-cauchy} wynika, że cały ciąg
$(M^n)$ zbiega do $M$ w normie $\mathbb S^2$.
Dla $s\le t$ mamy
$\E[M_t^n\mid\calF_s]=M_s^n$; zbieżność w $L^2$
i kontrakcja warunkowej wartości oczekiwanej pozwalają
przejść do granicy i otrzymać $\E[M_t\mid\calF_s]=M_s$.
Tym samym $M$ jest martyngałem.

Gdybyśmy wybrali inny ciąg aproksymujący $K^n\to H$,
to \eqref{eq:ito-doob-cauchy} dla różnicy pokazuje,
że oba procesy graniczne mają zerową normę odległości
w $\mathbb S^2$. Są więc nierozróżnialne.
Izometrię \eqref{eq:ito-full-iso} otrzymujemy,
przechodząc do granicy w odpowiedniej równości dla
procesów prostych, najpierw dla kwadratów, a następnie
stosując polaryzację. Dokładnie to samo daje
\eqref{eq:ito-process-limit}.
\end{proof}

\begin{wniosek}[Oszacowanie maksymalne]
Dla każdego $H\in\calH_T^2$ zachodzi
\[
\E\sup_{0\le t\le T}\left|\int_0^tH_s\,\dd W_s\right|^2
\le4\E\int_0^TH_s^2\,\dd s.
\tag{2.9}\label{eq:ito-max}
\]
\end{wniosek}
\begin{proof}
Jest to nierówność maksymalna Dooba dla martyngału
z twierdzenia \ref{th:ito-extension} wraz z izometrią.
\end{proof}

\begin{uwaga}[Co daje oszacowanie maksymalne?]
Sama izometria w chwili $T$ kontroluje tylko
$\E|I_T(H)|^2$. Oszacowanie \eqref{eq:ito-max}
kontroluje \emph{całą ścieżkę} całki stochastycznej.
Pozwala więc uzyskać zbieżność jednostajną na zwartym
przedziale w prawdopodobieństwie, a po wyborze podciągu
również zbieżność jednostajną prawie na pewno. To właśnie
uzasadnia wybór ciągłej wersji procesu całkowego.
\end{uwaga}


\begin{lemat}[Ciągły proces adaptowany jest przewidywalny]
Jeżeli $H$ jest rzeczywistym, adaptowanym procesem
o ciągłych trajektoriach, to jest przewidywalny.
Jeśli ponadto $\E\int_0^TH_t^2\,\dd t<\infty$,
to należy do $\calH_T^2$.
\end{lemat}
\begin{proof}
Dla każdego $n$ wybierzmy deterministyczny podział
$0=t_0^{(n)}<\cdots<t_{k_n}^{(n)}=T$ o średnicy
dążącej do zera i zdefiniujmy
$H_t^{(n)}=\sum_k H_{t_k^{(n)}}
\ind{(t_k^{(n)},t_{k+1}^{(n)}]}(t)$,
a w punkcie zero przyjmijmy $H_0^{(n)}=H_0$.
Każdy $H^{(n)}$ jest przewidywalny, bo
$H_{t_k^{(n)}}$ jest $\calF_{t_k^{(n)}}$-mierzalny.
Dla każdej ciągłej trajektorii mamy
$H_t^{(n)}\to H_t$ dla każdego $t$ dzięki
ciągłości w czasie. Granica punktowa funkcji
mierzalnych względem $\calP_T$ pozostaje
$\calP_T$-mierzalna. Druga część wynika wprost
z definicji przestrzeni $\calH_T^2$.
\end{proof}

\begin{przyklad}[Aproksymacja $H_t=g(W_t)$ i jej dokładność]
Niech $g:\R\to\R$ będzie funkcją globalnie Lipschitzowską
z pewną stałą $L\ge0$. Proces $H_t=g(W_t)$ jest
adaptowany i ciągły, a także należy do $\calH_T^2$,
ponieważ $|g(x)|\le|g(0)|+L|x|$ i
$\E W_t^2=t$. Na równomiernej siatce $t_k=kT/n$
przybliżmy $H$ przez
$H_t^{(n)}=g(W_{t_k})$ dla $t\in(t_k,t_{k+1}]$.
Wówczas
\[
\E\int_0^T|H_t^{(n)}-H_t|^2\,\dd t
\le L^2\sum_{k=0}^{n-1}\int_{t_k}^{t_{k+1}}
\E|W_t-W_{t_k}|^2\,\dd t
=\frac{L^2T^2}{2n}.
\]
Izometria Itô daje błąd średniokwadratowy całki
w chwili $T$ nie większy niż $L^2T^2/(2n)$,
a nierówność Dooba daje analogiczną kontrolę
całego procesu, z czterokrotnie większą stałą.
Takie oszacowanie jest punktem wyjścia
numerycznej aproksymacji całek stochastycznych.
\end{przyklad}

\section{Podstawowe własności wynikające z konstrukcji}
\begin{twierdzenie}[Lokalność w czasie i zależność od informacji]
\label{th:ito-properties}
Niech $H,K\in\calH_T^2$, $0\le s\le t\le T$.
Wówczas:
\begin{enumerate}[label=\textup{(\roman*)}]
\item $\displaystyle\int_0^tH_u\,\dd W_u
=\int_0^T H_u\ind{(0,t]}(u)\,\dd W_u$;
\item jeśli $H=K$ $\dd t\otimes\Prob$-prawie wszędzie
na $[0,t]\times\Omega$, to $I_u(H)=I_u(K)$ dla wszystkich
$u\le t$ poza jednym zbiorem zerowym;
\item $\displaystyle I_t(H)-I_s(H)=\int_0^TH_u\ind{(s,t]}(u)\,\dd W_u$;
\item jeśli $A\in\calF_s$, to
$\displaystyle\ind A\big(I_t(H)-I_s(H)\big)
=\int_s^t\ind A H_u\,\dd W_u$;
\item zachodzi warunkowa izometria
\[
\E\big[(I_t(H)-I_s(H))^2\mid\calF_s\big]
=\E\left[\left.\int_s^tH_u^2\,\dd u\,\right|\calF_s\right].
\]
\end{enumerate}
\end{twierdzenie}
\begin{proof}
Dla procesów prostych (i) i (iii) wynikają z samej
definicji sum. Dla ogólnych procesów przechodzimy
do granicy w $\calH_T^2$; pomnożenie procesu przez
wskaźnik deterministycznego przedziału nie zwiększa
normy. W (ii) izometria stosowana do $(H-K)\ind{(0,t]}$
daje równość całek w każdej chwili $u\le t$;
ciągłość obu procesów rozszerza ją z chwil wymiernych
na wszystkie chwile poza jednym zbiorem zerowym.

W (iv) proces $\ind A H_u\ind{(s,t]}(u)$ jest
przewidywalny, ponieważ $A\in\calF_s$ i na
$(s,t]$ informacja $A$ jest już dostępna. Dla procesu
prostego wzór wynika z przeniesienia wskaźnika $\ind A$
przed skończoną sumę. Dla ogólnego $H$ aproksymujemy
przez procesy proste i korzystamy z izometrii,
gdyż mnożenie przez $\ind A$ jest kontrakcją
w $\calH_T^2$.

W (v) dla procesu prostego mamy \eqref{eq:ito-conditional-var}.
Dla $H^n\to H$ w $\calH_T^2$ przyrosty całek
zbiegają w $L^2$, więc ich kwadraty zbiegają w $L^1$:
$\E|X_n^2-X^2|\le\norm{X_n-X}_2(\norm{X_n}_2+\norm X_2)$.
Jednocześnie $\int_s^t(H_u^n)^2\,\dd u$ zbiega
w $L^1(\Omega)$ do $\int_s^tH_u^2\,\dd u$,
co wynika z tej samej nierówności Cauchy'ego--Schwarza
zastosowanej w przestrzeni produktowej.
Ponieważ warunkowa wartość oczekiwana jest kontrakcją
w $L^1$, przechodzimy do granicy w tożsamości
warunkowej dla procesów prostych.
\end{proof}

\begin{twierdzenie}[Kwadrat całki i przewidywalny nawias]
\label{th:ito-bracket}
Dla $H\in\calH_T^2$ proces
\[
Q_t:=\left(\int_0^tH_s\,\dd W_s\right)^2
-\int_0^tH_s^2\,\dd s,\qquad 0\le t\le T,
\tag{2.10}\label{eq:ito-square-mg}
\]
jest martyngałem. Jeżeli przez $\langle M\rangle$
oznaczamy przewidywalny nawias ciągłego martyngału
kwadratowo całkowalnego z tomu II, to dla
$M_t=\int_0^tH_s\,\dd W_s$ mamy
$\langle M\rangle_t=\int_0^tH_s^2\,\dd s$.
W szczególności dla ciągłego martyngału
$[M]=\langle M\rangle$, więc
$[M]_t=\int_0^tH_s^2\,\dd s$.
\end{twierdzenie}
\begin{proof}
Własność martyngałową $M$ oraz warunkową izometrię
ustaliliśmy powyżej. Dla $s\le t$ zapisujemy
$M_t=M_s+(M_t-M_s)$ i rozwijamy kwadrat.
Po warunkowaniu względem $\calF_s$ składnik liniowy
$2M_s\E[M_t-M_s\mid\calF_s]$ znika. Pozostaje
\[
\E[M_t^2\mid\calF_s]
=M_s^2+\E\left[\left.\int_s^tH_u^2\,\dd u\,\right|\calF_s\right].
\]
Odjęcie $\E[\int_0^tH_u^2\,\dd u\mid\calF_s]$
daje $\E[Q_t\mid\calF_s]=Q_s$. Wszystkie wyrazy
są całkowalne na mocy izometrii.
Proces $A_t=\int_0^tH_s^2\,\dd s$ jest adaptowany,
ciągły i niemalejący, a więc przewidywalny.
Z definicji i jednoznaczności przewidywalnego nawiasu
ciągłego martyngału kwadratowo całkowalnego wynika
$\langle M\rangle=A$. W tomie II wykazaliśmy,
że dla ciągłych lokalnych martyngałów wariacja
kwadratowa i nawias przewidywalny są tym samym
ciągłym procesem; stąd ostatnia równość.
\end{proof}

\begin{uwaga}[Nie jest to jeszcze pełny wzór Itô]
Tożsamość \eqref{eq:ito-square-mg} identyfikuje
martyngałowy kompensator kwadratu. W kolejnym rozdziale
wyprowadzimy wzór dla ogólnego $F(t,X_t)$, a także
rachunek różniczkowy z poprawką wariacji kwadratowej.
Nie powinniśmy traktować symboli $(\dd W_t)^2=\dd t$
jako zwykłej algebry nieskończenie małych wielkości;
są one skrótem dla ścisłych twierdzeń o granicach sum.
\end{uwaga}


\begin{przyklad}[Martyngał bez niezależnych przyrostów]
Proces $M_t:=\int_0^tW_u\,\dd W_u$
jest martyngałem, ale jego przyrosty nie są niezależne
od przeszłości. Z warunkowej izometrii dla $s<t$
wynika
\[
\E[(M_t-M_s)^2\mid\calF_s]
=\E\left[\left.\int_s^tW_u^2\,\dd u\,\right|\calF_s\right]
=\int_s^t\big(W_s^2+(u-s)\big)\,\dd u
=(t-s)W_s^2+\frac{(t-s)^2}{2}.
\]
W drugim przejściu użyliśmy rozkładu
$W_u=W_s+(W_u-W_s)$ oraz niezależności
przyrostu od $\calF_s$. Prawa strona jest
rzeczywiście losowa dla $s>0$. Gdyby
$M_t-M_s$ był niezależny od $\calF_s$, jego
warunkowy drugi moment byłby stałą, co przeczy
otrzymanemu wzorowi. Ta obserwacja odróżnia
martyngałowość od własności niezależnych przyrostów.
\end{przyklad}

\section{Zatrzymywanie i całka stochastyczna}
Przy zatrzymywaniu strategii na losowym horyzoncie nie
wystarcza zwykłe pomnożenie przez wskaźnik zdarzenia
$\{\tau\le t\}$. Potrzebujemy procesu, który w chwili $u$
wskazuje, czy strategia jeszcze działa. Właśnie taki proces
jest przewidywalny.

\begin{twierdzenie}[Zatrzymanie całki Itô]\label{th:ito-stopping}
Niech $H\in\calH_T^2$ i niech $\tau$ będzie czasem
zatrzymania o wartościach w $[0,T]$. Wówczas proces
$H_u\ind{(0,\tau]}(u)$ jest przewidywalny i należy
do $\calH_T^2$, a dla wszystkich $t\in[0,T]$
zachodzi, poza jednym zbiorem zerowym,
\[
\int_0^{t\wedge\tau}H_u\,\dd W_u
=\int_0^t H_u\ind{\{u\le\tau\}}\,\dd W_u.
\tag{2.11}\label{eq:stopped-ito}
\]
\end{twierdzenie}
\begin{proof}
Proces $u\mapsto\ind{\{u\le\tau\}}$ jest adaptowany
i lewostronnie ciągły dla każdej trajektorii na $(0,T]$,
więc jest przewidywalny. Iloczyn z przewidywalnym $H$
też jest przewidywalny, a jego norma w $\calH_T^2$
nie przekracza normy $H$.

Najpierw niech $H$ będzie procesem prostym. Ustalmy
ciąg czasów $\tau_n:=2^{-n}\lceil2^n\tau\rceil\wedge T$,
gdzie w razie potrzeby ostatni przedział siatki skracamy
do $T$. Każdy $\tau_n$ jest czasem zatrzymania o skończenie
wielu wartościach, a $\tau_n\downarrow\tau$.
Po wspólnym rozdrobnieniu siatki procesu $H$ i siatki
czasu $\tau_n$ można zapisać proces
$H\ind{(0,\tau_n]}$ jako prosty przewidywalny:
współczynnik na odcinku $(u,v]$ zawiera wskaźnik
$\{\tau_n\ge v\}=\{\tau_n>u\}$, który jest
$\calF_u$-mierzalny. Z definicji sum wynika dokładnie
$I_t(H\ind{(0,\tau_n]})=I_{t\wedge\tau_n}(H)$.
Ponieważ $H$ jest kwadratowo całkowalny,
\[
\E\int_0^T H_u^2
|\ind{\{u\le\tau_n\}}-\ind{\{u\le\tau\}}|^2\,\dd u
\longrightarrow0
\]
na mocy zbieżności zdominowanej względem
$\dd t\otimes\dd\Prob$; pojedyncza chwila $u=\tau$
nie wpływa na całkę po czasie. Izometria Itô pozwala
przejść do granicy po prawej, a ciągłość trajektorii
$I(H)$ po lewej. Dostajemy tezę dla prostego $H$.

W przypadku ogólnym wybierzmy proste $H^m\to H$
w $\calH_T^2$. Po zatrzymaniu zachodzi ta sama
zbieżność w $\calH_T^2$, gdyż mnożenie przez
$\ind{(0,\tau]}$ jest kontrakcją.
Nierówność maksymalna \eqref{eq:ito-max} daje zbieżność
całych procesów $I(H^m)\to I(H)$ oraz
$I(H^m\ind{(0,\tau]})\to I(H\ind{(0,\tau]})$
w normie $\mathbb S^2$. Tożsamość dla $H^m$
przechodzi zatem do granicy. Równość w chwilach wymiernych
i ciągłość procesów zapewniają wspólny zbiór pełnego
prawdopodobieństwa, na którym zachodzi dla wszystkich $t$.
\end{proof}

\begin{zadanie}[Zatrzymany ruch Browna]
Niech $a>0$, a proces Wienera będzie określony na
$[0,\infty)$. Niech $\tau_a=\inf\{t\ge0:W_t=a\}$.
Dla ustalonego $T>0$
wyznaczyć $\int_0^T\ind{\{t\le\tau_a\}}\,\dd W_t$
i jego wartość oczekiwaną.

\emph{Rozwiązanie.} Z twierdzenia \ref{th:ito-stopping}
i równości $\int_0^t1\,\dd W_s=W_t$ otrzymujemy
$\int_0^T\ind{\{t\le\tau_a\}}\,\dd W_t=W_{T\wedge\tau_a}$.
Ponieważ czas $T\wedge\tau_a$ jest ograniczony,
próbkowanie opcjonalne lub bezpośrednia martyngałowość
całki daje $\E W_{T\wedge\tau_a}=0$.
Nie wolno bez dodatkowych warunków przechodzić tu do
$T\to\infty$ pod wartością oczekiwaną:
$W_{\tau_a}=a$ prawie na pewno, ale $\E W_{\tau_a}=a$,
a nie zero. Przy granicy traci się odpowiednią
jednostajną całkowalność.
\end{zadanie}

\section{Lokalna całkowalność: wyjście poza przestrzeń kwadratowalną}
Warunek $\E\int_0^T H_t^2\,\dd t<\infty$ jest
wygodny do konstrukcji, ale w wielu zastosowaniach
okazuje się zbyt mocny. Proces może mieć dobrze
zdefiniowaną całkę mimo nieskończonej wartości oczekiwanej
energii, jeśli energia każdej pojedynczej trajektorii
jest skończona.

\begin{definicja}[Całkowalność lokalna względem procesu Wienera]
Mówimy, że przewidywalny proces $H$ należy do
$\calH_{\mathrm{loc},T}^2$, gdy
$\int_0^T H_t^2\,\dd t<\infty$ prawie na pewno.
Nie zakładamy skończoności wartości oczekiwanej tej całki.
\end{definicja}

\begin{twierdzenie}[Konstrukcja całki Itô przez lokalizację]\label{th:ito-local}
Dla każdego $H\in\calH_{\mathrm{loc},T}^2$
istnieje jednoznaczny, z dokładnością do nierozróżnialności,
ciągły lokalny martyngał $I(H)$, taki że dla
$\tau_n:=\inf\{t\in[0,T]:\int_0^t H_s^2\,\dd s\ge n\}\wedge T$
zachodzi
\[
I_{t\wedge\tau_n}(H)
=\int_0^tH_s\ind{\{s\le\tau_n\}}\,\dd W_s,
\qquad 0\le t\le T.
\tag{2.12}
\]
Procesy po prawej są zdefiniowane przez
poprzednią konstrukcję w $L^2$, a $\tau_n\uparrow T$
prawie na pewno. Całka zachowuje liniowość i własność
zatrzymywania.
\end{twierdzenie}
\begin{proof}
Proces $A_t=\int_0^tH_s^2\,\dd s$ jest adaptowany,
ciągły i niemalejący, więc $\tau_n$ jest czasem
zatrzymania. Ciągłość $A$ zapewnia
$A_{\tau_n}\le n$, chyba że $\tau_n=T$ i $A_T<n$;
w obu przypadkach nierówność pozostaje prawdziwa.
Zatem
\[
\E\int_0^T H_s^2\ind{\{s\le\tau_n\}}\,\dd s
=\E A_{\tau_n}\le n.
\]
W konsekwencji każdy proces
$H^{(n)}:=H\ind{(0,\tau_n]}$ należy do $\calH_T^2$.
Dla $m\ge n$ mamy
$H^{(m)}\ind{(0,\tau_n]}=H^{(n)}$ niemal wszędzie.
Twierdzenie \ref{th:ito-stopping} daje więc
$I_{t\wedge\tau_n}(H^{(m)})=I_t(H^{(n)})$ dla wszystkich
$t$ poza wspólnym zbiorem zerowym.
Można zatem skleić ciągłe procesy $I(H^{(n)})$
w jedną ciągłą trajektorię: definiujemy ją w chwili $t$
przez $I_t(H^{(n)})$, jeśli $t\le\tau_n$, co jest
niezależne od dostatecznie dużego $n$.
Ponieważ $A_T<\infty$ prawie na pewno, dla każdego
$\omega$ poza zbiorem zerowym przy dostatecznie dużym
$n$ mamy $\tau_n(\omega)=T$; sklejanie obejmuje więc
cały przedział. Otrzymany proces jest lokalnym martyngałem,
bo każdy zatrzymany proces $I_{\cdot\wedge\tau_n}(H)$
jest martyngałem kwadratowo całkowalnym.
Jednoznaczność wynika z jednoznaczności wszystkich
zatrzymanych procesów. Liniowość otrzymujemy,
lokalizując równocześnie dwa integrandy przez
wspólną sekwencję czasów opartą na sumie ich energii.
\end{proof}

\begin{uwaga}[Kres własności całki]
Dla lokalnie kwadratowo całkowalnego procesu $H$ całka
$I_t(H)$ \emph{nie musi} należeć do $L^2(\Omega)$ ani
być prawdziwym martyngałem. Jej podstawową własnością
jest lokalna martyngałowość. Nie wolno więc bezwarunkowo
stosować wzoru $\E I_t(H)=0$ ani izometrii w postaci
z nieskończoną prawą stroną. Poprawny jest zawsze
wzór lokalny, po zatrzymaniu na $\tau_n$.
\end{uwaga}

\begin{przyklad}[Proces lokalnie całkowalny, lecz nie w $\calH_T^2$]
Niech $Z$ będzie zmienną losową niezależną od $W$,
$Z<\infty$ prawie na pewno, ale $\E Z^2=\infty$.
Rozszerzmy filtrację przez ujawnienie $Z$ w chwili $0$
i zdefiniujmy $H_t=Z$. Proces jest przewidywalny i
$\int_0^TH_t^2\,\dd t=TZ^2<\infty$ prawie na pewno,
lecz jego wartość oczekiwana jest nieskończona.
Całka lokalna ma postać $I_t(H)=ZW_t$.
Proces zatrzymujemy po przekroczeniu poziomu energii;
w tej prostej sytuacji można też bezpośrednio lokalizować
przez zdarzenia $\{|Z|\le n\}\in\calF_0$.
\end{przyklad}

\section{Obliczanie całek: trzy szczegółowe przykłady}
\begin{zadanie}[Podstawowa całka $\int W\,\dd W$]\label{zad:W-dW}
Obliczyć $\int_0^T W_t\,\dd W_t$ bez używania wzoru Itô.

\emph{Rozwiązanie.}
Proces $H_t=W_t$ jest adaptowany i ciągły,
a więc przewidywalny. Ponadto
$\E\int_0^TW_t^2\,\dd t=\int_0^Tt\,\dd t=T^2/2<\infty$.
Na równomiernej siatce $t_k=kT/n$ rozważmy proces
$H_t^{(n)}=W_{t_k}$ dla $t\in(t_k,t_{k+1}]$.
Z niezależności przyrostów
\[
\E\int_0^T|H_t^{(n)}-W_t|^2\,\dd t
=\sum_k\int_{t_k}^{t_{k+1}}(t-t_k)\,\dd t
=\frac{T^2}{2n}\longrightarrow0.
\]
Izometria Itô daje więc zbieżność sum
$\sum_kW_{t_k}\Delta_kW$ do szukanej całki w $L^2$.
Z tożsamości algebraicznej
\[
2\sum_kW_{t_k}\Delta_kW
=W_T^2-\sum_k(\Delta_kW)^2
\]
oraz zbieżności wariacji kwadratowej z tomu II
$\sum_k(\Delta_kW)^2\to T$ w $L^2$ otrzymujemy
\[
\int_0^TW_t\,\dd W_t=\frac12(W_T^2-T)
\quad\text{prawie na pewno}.
\tag{2.13}\label{eq:W-dW}
\]
Warto zauważyć, że średnia prawej strony wynosi zero,
a jej wariancja jest równa
$\frac14\Var(W_T^2)=T^2/2$, zgodnie z izometrią.
\end{zadanie}

\begin{zadanie}[Całka $\int W^2\,\dd W$ bez wzoru Itô]
Obliczyć $\int_0^T W_t^2\,\dd W_t$ bez stosowania
rachunku Itô do funkcji $x\mapsto x^3$.

\emph{Rozwiązanie.} Proces $W_t^2$ jest przewidywalny,
a $\E\int_0^T W_t^4\,\dd t
=3\int_0^T t^2\,\dd t=T^3<\infty$.
Dla siatki równomiernej $t_k=k\Delta$, $\Delta=T/n$,
przyrostów $\Delta_kW=W_{t_{k+1}}-W_{t_k}$ mamy
dokładną tożsamość algebraiczną
\[
W_T^3
=3\sum_kW_{t_k}^2\Delta_kW
+3\sum_kW_{t_k}(\Delta_kW)^2
+\sum_k(\Delta_kW)^3.
\tag{2.14}\label{eq:ito-cube-expand}
\]
Najpierw udowodnimy zbieżność pierwszej sumy do całki Itô.
Na przedziale $[t_k,t_{k+1}]$, pisząc
$W_t=W_{t_k}+\delta_t$ z przyrostem niezależnym od
$\calF_{t_k}$, otrzymujemy
$\E(W_t^2-W_{t_k}^2)^2
=4t_k(t-t_k)+3(t-t_k)^2$.
Wobec tego
\[
\E\int_0^T|W_t^2-W_{t_k(t)}^2|^2\,\dd t
\le 4T^2\Delta+3T\Delta^2\longrightarrow0,
\]
gdzie $t_k(t)$ oznacza lewy koniec przedziału
zawierającego $t$. Izometria daje pożądaną zbieżność
pierwszej sumy w $L^2$.

Drugi składnik rozdzielamy na
\[
\sum_kW_{t_k}(\Delta_kW)^2
=\sum_kW_{t_k}\Delta+R_n,\qquad
R_n:=\sum_kW_{t_k}[(\Delta_kW)^2-\Delta].
\]
Poszczególne składniki $R_n$ są różnicami
martyngałowymi, więc są ortogonalne w $L^2$. Ponieważ
$\E[(\Delta_kW)^2-\Delta]^2=2\Delta^2$,
\[
\E R_n^2
=2\Delta^2\sum_k\E W_{t_k}^2
\le2T^2\Delta\longrightarrow0.
\]
Suma Riemanna $\sum_kW_{t_k}\Delta$ dąży do
$\int_0^TW_t\,\dd t$ w $L^2$ (na przykład
z nierówności Cauchy'ego--Schwarza oraz ciągłości
średniokwadratowej $W$). Wreszcie składniki
$(\Delta_kW)^3$ są niezależne, mają średnią zero
i szóste momenty $15\Delta^3$. Stąd
$\E|\sum_k(\Delta_kW)^3|^2=15n\Delta^3=15T\Delta^2\to0$.
Przechodząc w \eqref{eq:ito-cube-expand} do granicy w $L^2$
otrzymujemy
\[
\int_0^TW_t^2\,\dd W_t
=\frac13W_T^3-\int_0^TW_t\,\dd t.
\tag{2.15}\label{eq:W2-dW}
\]
Ta tożsamość jest zapowiedzią wzoru Itô dla
funkcji $x^3$ i pokazuje, skąd bierze się dodatkowy
wyraz względem zwykłego rachunku różniczkowego.
\end{zadanie}

\begin{zadanie}[Strategia zależna od pierwszej obserwacji]
Załóżmy, że inwestor obserwuje $W_{T/2}$, a następnie
w drugiej połowie przedziału utrzymuje pozycję równą
$1$, jeśli $W_{T/2}>0$, i zero w przeciwnym przypadku.
Wyznaczyć prawo końcowego zysku
$X=\int_0^T\ind{\{W_{T/2}>0\}}\ind{(T/2,T]}(t)\,\dd W_t$.

\emph{Rozwiązanie.} Funkcja podcałkowa jest prostym
procesem przewidywalnym: decyzja jest
$\calF_{T/2}$-mierzalna. Z definicji
$X=\ind{\{W_{T/2}>0\}}(W_T-W_{T/2})$.
Przyrost $Y=W_T-W_{T/2}\sim\Normal(0,T/2)$ jest
niezależny od $W_{T/2}$. Z symetrii rozkładu normalnego
$\Prob(W_{T/2}>0)=1/2$. Zatem rozkład $X$ jest mieszaniną
$\frac12\delta_0+\frac12\Normal(0,T/2)$.
Nie jest rozkładem normalnym, ponieważ ma atom
w zerze o masie $1/2$. Jednocześnie
$\E X=0$ i $\Var(X)=(1/2)(T/2)=T/4$,
co potwierdza izometria Itô.
\end{zadanie}

\section{Przewidywalność nie jest kosmetycznym założeniem}
\begin{przyklad}[Integrand antycypujący]
Przypuśćmy błędnie, że moglibyśmy w całej konstrukcji
brać losowy współczynnik $\xi=W_T$ niezależnie od
czasu. Formalna suma dawałaby
$\sum_kW_T\Delta_kW=W_T^2$.
Jeśli nierozważnie użyć izometrii Itô, otrzymalibyśmy
$\E W_T^4\overset{?}=\E\int_0^TW_T^2\,\dd t=T^2$.
Lewa strona wynosi jednak $3T^2$.
Sprzeczność bierze się stąd, że $W_T$ nie jest
$\calF_{t_k}$-mierzalny dla $t_k<T$.
Nie oznacza to, że nigdy nie można całkować
antycypująco: istnieje na przykład całka Skorochoda,
lecz ma inną konstrukcję i inne własności.
\end{przyklad}

\begin{przyklad}[Dlaczego całka Itô nie musi mieć rozkładu normalnego]
Niech $Z\in\{0,1\}$ będzie zmienną Bernoulliego
z parametrem $1/2$, niezależną od procesu Wienera,
i niech $Z$ będzie znana w chwili $0$.
Wtedy $H_t=Z$ jest przewidywalny, zaś
$\int_0^TH_t\,\dd W_t=ZW_T$.
Prawo tej całki jest mieszaniną
$\frac12\delta_0+\frac12\Normal(0,T)$.
Choć każdy składnik sumy w definicji całki ma
normalny przyrost warunkowo względem przeszłości,
całka z \emph{losowym} integrandem nie musi być
normalna. Teza o gaussowskości całki z rozdziału
pierwszego obowiązywała dla funkcji deterministycznych.
\end{przyklad}


% Uzupełnienie audytowe rozdzialu 2: osobne uzasadnienia dwoch przejsc granicznych.
\section{Dwa kryteria przybliżania całek i ich ograniczenia}
\begin{lemat}[Lewostronna aproksymacja procesu ciągłego]
\label{lem:step-approx-ito}
Niech $H$ będzie ciągłym, adaptowanym procesem na
$[0,T]$, a $\sigma$ przewidywalnym procesem spełniającym
$\int_0^T\sigma_s^2\,\dd s<\infty$ prawie na pewno.
Niech $\pi_n$ będą deterministycznymi podziałami
$[0,T]$ o średnicy dążącej do zera, a
$H_s^{(n)}=H_{t_k}$ dla $s\in(t_k,t_{k+1}]$.
Wówczas, w sensie całek zlokalizowanych,
\[
\sup_{t\le T}\left|
\int_0^t(H_s^{(n)}-H_s)\sigma_s\,\dd W_s\right|
\longrightarrow0\quad\text{według prawdopodobieństwa}.
\]
W szczególności sumy lewostronne
$\sum_kH_{t_k}(W_{t_{k+1}\wedge t}-W_{t_k\wedge t})$
zbiegają do $\int_0^tH_s\,\dd W_s$ jednostajnie
według prawdopodobieństwa.
\end{lemat}
\begin{proof}
Ponieważ każda trajektoria $H$ jest jednostajnie
ciągła na $[0,T]$, mamy
$\norm{H^{(n)}-H}_\infty\to0$ prawie na pewno.
Wprowadźmy czas zatrzymania
$\tau_m=\inf\{t:\int_0^t\sigma_s^2\,\dd s\ge m\}
\wedge T$ oraz $\rho_m=\inf\{t:|H_t|\ge m\}\wedge T$.
Rozważmy $\theta_m=\tau_m\wedge\rho_m$.
Na $[0,\theta_m]$ mamy
$\int_0^{\theta_m}\sigma_s^2\,\dd s\le m$ oraz
$|H_s|\le m$; dla procesów $H^{(n)}$ po przekroczeniu
$\theta_m$ nie całkujemy, a ich wartości na
$[0,\theta_m]$ są ograniczone przez $m$.
Zatem na mocy izometrii i zbieżności zdominowanej
na przestrzeni $\Omega$,
\[
\E\int_0^{\theta_m}|H_s^{(n)}-H_s|^2
\sigma_s^2\,\dd s
\le \E\left[
\Big(\sup_{s\le\theta_m}|H_s^{(n)}-H_s|\Big)^2
\int_0^{\theta_m}\sigma_s^2\,\dd s\right]\longrightarrow0,
\]
ponieważ supremum na zatrzymanym przedziale jest
ograniczone przez $2m$ i zbiega do zera prawie na pewno.
Nierówność Dooba zapewnia
zbieżność zatrzymanych całek w $L^2$ normy supremum.
Ponieważ $\Prob(\theta_m<T)\to0$, usunięcie
lokalizacji daje tezę. Dla $\sigma\equiv1$
pierwsza równość wprost identyfikuje granicę
sum lewostronnych.
\end{proof}

\begin{zadanie}[Dlaczego zbieżność w $L^2$ nie daje zbieżności
ścieżkowej dla całego ciągu?]
Niech $(Y_n)$ będzie ciągiem niezależnych zdarzeń
z $\Prob(Y_n=1)=1/n$ i $\Prob(Y_n=0)=1-1/n$
($n\ge2$). Wyznaczyć rodzaje zbieżności $Y_n$ do zera.

\emph{Rozwiązanie.} Mamy $\E|Y_n|^2=1/n\to0$,
a zatem $Y_n\to0$ w $L^2$ i według prawdopodobieństwa.
Ponieważ $\sum_n\Prob(Y_n=1)=\infty$ i zdarzenia są
niezależne, drugi lemat Borela--Cantellego daje
$\Prob(Y_n=1\ \text{nieskończenie często})=1$.
Zatem nie ma zbieżności do zera prawie na pewno.
Wyjaśnia to, dlaczego w konstrukcji całki Itô potrzebny
był wybór odpowiedniego podciągu do uzyskania ciągłej
wersji procesu: z samej zbieżności w $L^2$ całego
ciągu nie wolno wyciągać wniosku o zbieżności
ścieżkowej całego ciągu.
\end{zadanie}


\section{Co udowodniliśmy, a co pozostaje do zrobienia?}
Całka Itô jest obecnie określona dla wszystkich
procesów przewidywalnych o skończonej oczekiwanej energii,
a przez lokalizację również dla procesów o energii
skończonej niemal na pewno. Udowodniliśmy izometrię,
martyngałowość, ciągłość, zatrzymywanie i identyfikację
wariacji kwadratowej. Na tej podstawie można już
formułować stochastyczne równania całkowe
$X_t=X_0+\int_0^tb_s\,\dd s+\int_0^t\sigma_s\,\dd W_s$.
Nie udowodniliśmy jeszcze ogólnego wzoru Itô,
nie skonstruowaliśmy całki Stratonovicha i nie badaliśmy
istnienia oraz jednoznaczności rozwiązań SDE.
Te zagadnienia będą przedmiotem kolejnych rozdziałów.

\begin{uwaga}[Konwencje na dalszą część tomu]
Przez $\int_0^tH_s\,\dd W_s$ będziemy rozumieli
całkę Itô, o ile wyraźnie nie zaznaczymy innej
konwencji. Skrót $\dd X_t=b_t\,\dd t+\sigma_t\,\dd W_t$
zawsze będziemy interpretować przez odpowiednie
równanie całkowe. Całki po trajektoriach z miarą
Wienera i całki względem pojedynczego procesu
pozostają odrębnymi konstrukcjami, nawet gdy
w szczególnych przykładach łączy je wzór Feynmana--Kaca.
\end{uwaga}
% ============================================================================
% ROZDZIAL 3 — WZOR ITO
% ============================================================================
